\section{From Diagrams to Asymptotic GFOM Dynamics}
\label{sec:dynamics}


The traffic distribution captures the limiting behavior of all scalar-valued,
permutation-invariant polynomials. In this section, we show how to leverage this
information to derive the limiting empirical laws of vector-valued,
permutation-invariant polynomials. Our main application is a description of the
limiting dynamics of GFOM.


We will mostly work under the assumption that the input matrices $\bA = \bA^{(n)}$ satisfy the strong cactus property, which we recall is the statement that $\frac{1}{n}\E_\bA z_{\alpha}(\bA) \to 0$ as $n \to \infty$ for all non-cactus $\alpha$ (i.e., all $\alpha \in \cA \setminus \cC$, a statement about \emph{scalar} graph polynomials).
In~\Cref{sec:dynamics-reg} we will briefly suspend this assumption to discuss punctured matrices, so as to connect to the setting of~\Cref{sec:universality}.

We tackle two tasks in this section:
\begin{enumerate}
    \item First, we study the joint asymptotic limit of the empirical distributions of the vector diagrams $\bmz_\al(\bA)$ over $\al \in \calA_1$. 
    Assuming the strong cactus property, we show that only the small subset of \emph{treelike} $\al \in \cT_1$ are asymptotically nonzero in the $z$-basis, in a sense to be made precise below.
    We then show that the asymptotic algebra of the treelike diagrams is isomorphic to a \emph{Wick algebra}, an algebra defined by a family of Gaussian random variables.
    This will give a precise version of \Cref{prop:convergence-distribution}.
    \item Second, we work with the asymptotic limit of treelike diagrams to identify a generalized Onsager correction, derive a treelike Approximate Message Passing algorithm, and prove its state evolution over arbitrary input matrices having the strong cactus property and a limiting diagonal distribution.
    This will give a precise version of \Cref{thm:intro-onsager}.
\end{enumerate}

\subsection{Asymptotic limit of the vector diagrams}
\label{sec:vector-limit}

In this section, given a family $(X_i)_{i\in I}$ and
$J\subseteq I$, we will write as a shortcut $X_J = (X_j)_{j\in J}$. 

Recall that $\cC_1$ denotes the set of rooted cactuses and $\cT_1 \subseteq \cA_1$ denotes the set of rooted trees with hanging cactuses.
We call the diagrams in $\cT_1$ \emph{treelike}, and we call \emph{Gaussian trees} the subset of diagrams $\cG_1 \subseteq \cT_1$ such that the root has degree exactly $1$ after removing hanging cactuses.

\begin{definition}[Type]
    For each $\tau\in\cT_1$, let
    $\type(\tau) \in \N^{\cG_1\cup \calC_1}$, where $\type(\tau)_\alpha$ count the
    number of copies of $\alpha\in \calG_1\cup \calC_1$ attached to the root of $\tau$,
    with the additional convention that
    $\type(\tau)_\alpha=0$ for all
    $\alpha\in\calG_1$ that has cactuses hanging at the root.
\end{definition}

The following theorem identifies the limiting distribution of $\bz_{\calA_1}(\bA)$ under the strong cactus property.
We refer the reader to~\Cref{sec:convergence-process} for the definition of convergence in distribution for random elements indexed by countably infinite index sets.

\begin{theorem}
\label{thm:convergence-distribution}
    Assume that $\bA=\bA^{(n)}$ satisfies~\cref{eq:tightnessNorm}, has the strong cactus property, and a limiting diagonal distribution.
    Then,
    \[
        \samp(\bmz_{\cA_1}(\bA)) \tod Z_{\cA_1}^{\infty}\,,
    \]
    where $Z_{\calA_1}^{\infty} \in \R^{\calA_1}$ is a random variable satisfying the following properties:
    \begin{enumerate}
        \item $Z_{\alpha}^\infty = 0$ for all non-treelike $\alpha$.

        \item Conditioned on $Z^\infty_{\calC_1}$, $Z^\infty_{\calG_1}$ is a centered
        Gaussian process with covariance $\bm \Sigma^\infty$ from~\cref{eq:covariance}.
        
        \item Let $\He$ denote the Wick product (\Cref{def:wick}). Then
        for every $\tau\in\calT_1$,
        \[
            Z^\infty_\tau
            = \He_{\type(\tau)}(Z^\infty_{\calG_1}\,;\,\bm \Sigma^\infty)\cdot \prod_{\sigma\in \calC_1} (Z^\infty_{\sig})^{\type(\tau)_\sigma}\,.
        \]
    \end{enumerate}
\end{theorem}

\Cref{thm:convergence-distribution} shows how
the limiting algebra $Z^\infty_{\calA_1}$
of permutation-invariant, vector-valued polynomials in
$\bA$ can be derived from $Z^\infty_{\calC_1}$.
Although we have not specified the description of the law of $Z^\infty_{\calC_1}$, it is fully
determined by the limiting diagonal distribution of $\bA$. For example, when $\bA$ further
satisfies the
{\em factorizing} strong cactus property, $Z^\infty_{\calC_1}$ is deterministic:

\begin{proposition}
    \label{prop:factorizing-deterministic}
    If $\bA$ satisfies the factorizing strong 
    cactus property and~\cref{eq:tightnessNorm},
    then the conclusion of~\Cref{thm:convergence-distribution} holds with the additional property that for every
    $\sigma\in\calC_1$,
    \[
        Z^\infty_{\sigma} = \prod_{\rho\in\cyc(\sigma)} \left(\lim_{n\to\infty} \frac 1n \E z_\rho(\bA)\right)\,.
    \]
\end{proposition}

\begin{proof}
    Let $\sigma\in\calC_1$.
    The first moment of
    $\samp(\bz_\sigma(\bA))$ is
    \begin{align}
        \E \samp(\bz_\sigma(\bA)) = \frac 1n \sum_{i=1}^n \E_\bA \bz_\sigma(\bA)[i] = \frac 1n \E_\bA z_{\sigma_0}(\bA)\,,\label{eq:first-prop53}
    \end{align}
    where $\sigma_0$ is the unrooted version of $\sigma$. As $n\to\infty$,~\cref{eq:first-prop53} converges to the deterministic
    constant
    \[
        \kappa_{\sigma_0} \defeq \lim_{n\to\infty} \frac 1n \E_\bA z_{\sigma_0}(\bA) = \prod_{\rho\in\cyc(\sigma_0)} \left(\lim_{n\to\infty} \frac 1n \E z_\rho(\bA)\right)
    \]
    by the factorizing cactus property.
    
    We now switch to the second
    moment,
    \[
        \E \samp(\bz_\sigma(\bA))^2 =  \frac 1n \E_\bA \sum_{i=1}^n  \bz_\sigma(\bA)[i]^2\,.
    \]
    Expand the scalar polynomial $\sum_{i=1}^n \bz_\sigma(\bA)[i]^2$
    in the $z$-basis. The support of that
    expansion is
    the set of diagrams that can be
    obtained by grafting two copies
    of $\sigma$ at the root and merging
    pairs of vertices across the two
    different copies. By the strong cactus property, it suffices to find
    which cactuses can be obtained in
    this way. By~\Cref{lem:prod-cactus-cactus}, the only
    cactus that can occur in this way
    has no merging, and it contributes
    $\kappa_{\sig_0}^2$ by the factorizing cactus property.
    Thus,
    \[
        \E \samp(z_\sig(\bA))^2 = \kappa_{\sig_0}^2 + o(1) = \left(\E \samp(z_\sig(\bA)\right)^2 + o(1)\,.
    \]

    We showed that $\samp (z_\sigma(\bA))$
    converges to the desired 
    deterministic quantity in expectation,
    and furthermore that
    its variance converges to $0$.
    This implies that it converges to
    the constant in distribution. By
    unicity of the limit in distribution,
    $Z_\sigma^\infty$
    equals that constant almost surely.
\end{proof}

However, if we drop the factorizing cactus property assumption,
the variables $Z^\infty_{\calC_1}$ may no longer be deterministic. 
For example, this can be the case
when $\bA$ is a block-structured matrix as
in~\Cref{sec:wigner-block}:
 
 \begin{example}\label{ex:block-structured}
    Let $\bm A_1^{(n)}$ and $\bm A_2^{(n)}$
    be two $n\times n$ matrices
    satisfying the assumptions of~\Cref{thm:convergence-distribution}. Define
    the $2n\times 2n$ matrix,
    \[
        \bmA^{(2n)} = \begin{bmatrix} \bmA_1^{(n)} & \bm 0\\\bm 0 & \bmA_2^{(n)} \end{bmatrix}
    \]
    From the block-diagonal structure, for any $\alpha\in \calA_1$,
    \[
        \bmz_{\alpha}(\bmA)[i] = \begin{cases} \bmz_{\alpha}(\bmA_1)[i] & \text{if $i\in [n]$}\\
        \bmz_{\alpha}(\bmA_2)[i-n] & \text{if $i\in [2n]\setminus [n]$}\end{cases}
    \]
    Hence, the law of $Z^\infty_{\calC_1}(\bmA)$ is a
    uniform mixture of the law of $Z^\infty_{\calC_1}(\bmA_1)$
    and that of $Z^\infty_{\calC_1}(\bmA_2)$.
\end{example}

\noindent We will prove a generalization of~\Cref{ex:block-structured} later; see~\Cref{lem:block-matrix-cactus-limits}.

In~\Cref{ex:block-structured}, the randomness of $Z^\infty_{\calC_1}$
may be viewed as coming solely from the $\samp(\cdot)$ operator,
but this is not always the case.
For instance, our model also captures orthogonally invariant distributions
that do not satisfy the traffic concentration property:

\begin{example}\label{ex:finetti}
    Let $(\lambda_n)_{n\ge 1}$
    be an exchangeable sequence of random variables in $[-1,1]$ and consider
    \[\bA^{(n)} = (\bQ^{(n)})^\top \diag(\lambda_1, \ldots, \lambda_n) \bQ^{(n)}\,,\]
    for Haar-distributed matrices $\bQ^{(n)} \in O(n)$, independent from $(\lambda_n)$. By de Finetti's
    theorem, there exists a latent random
    probability measure $\mu$ almost
    surely supported on $[-1,1]$ such that
    conditionally on $\mu$, $\lambda_1,\lambda_2,\ldots$ are i.i.d.\ with common law $\mu$.
    By~\Cref{thm:moments}, $\bA^{(n)}$ satisfies the strong
    cactus property conditionally on
    $\mu$, so it also satisfies the strong
    cactus property unconditionally.

    Applying~\Cref{thm:convergence-distribution} and \Cref{prop:factorizing-deterministic},
    we get that conditionally on $\mu$,
    $\samp(z_{\calC_1}(\bA))$ converges in
    distribution to
    \begin{align}
        \left(\prod_{\rho\in\cyc(\sigma)} \kappa_{|\rho|}(\mu)\right)_{\sigma\in \calC_1}\,,\label{eq:random-cumulants}
    \end{align}
    where $(\kappa_q(\mu))_{q\ge 1}$
    are the free cumulants of $\mu$.
    Therefore, {unconditionally}, 
    $\samp(z_{\calC_1}(\bA))$ converges
    in distribution to the {\em random}
    quantity~\cref{eq:random-cumulants}.
\end{example}

\noindent Note that \Cref{ex:block-structured,ex:finetti} do not contradict~\Cref{prop:factorizing-deterministic} because in these
examples,
$\bA^{(n)}$
typically does not 
satisfy the factorizing
cactus property.

\subsubsection{Non-treelike diagrams are asymptotically negligible}
\label{sec:treelike}

The remainder of~\Cref{sec:vector-limit} is dedicated to
the proof of~\Cref{thm:convergence-distribution}. In the whole proof, we drop the dependence
of $z_\alpha$ and $w_\alpha$ on $\bA$ 
to lighten notation.
We start by proving that
non-treelike diagrams are negligible.

\begin{lemma}
    \label{lem:non-treelike}
    Suppose that $\bA$ satisfies the
    strong cactus property. Then for each non-treelike $\alpha$,
    \[
        \samp(\bz_\alpha)\overset{L^2}{\longrightarrow} 0\,.
    \]
\end{lemma}

\begin{proof}
    By definition, we have
    \begin{align*}
        \E \samp(\bz_\alpha)^2 &= \frac 1n \sum_{i=1}^n \E \left[\left(\bz_\alpha^{ 2}\right)[i]\right]
    \intertext{By~\Cref{claim:only-trees}, we can expand $\bz_\alpha^{ 2}$ in the
    $\bz$-basis to obtain, for some constant coefficients
    $c_\beta$}
        &=  \frac 1 n \sum_{i=1}^n  \sum_{\beta\in \calA_1\setminus \calT_1} c_\beta \E \bz_\beta[i]\\
        &= \frac 1n \sum_{\beta\in \calA_0\setminus \calT_0} c_\beta' \E z_\beta
    \end{align*}
    for some other constant coefficients $c_\beta'$. Since no diagram in
    $\calA_0\setminus \calT_0$ is a cactus, by the strong
    cactus property, we get $\E \samp(\bz_\alpha)^2\underset{n\to\infty}\longrightarrow 0$, as desired.
\end{proof}



\subsubsection{Asymptotic limit of the treelike diagrams}
\label{sec:asymptotic-space}

Next, we analyze the treelike diagrams. All results in~\Cref{sec:asymptotic-space} are purely
combinatorial, meaning that they hold for arbitrary
$\bA\in\R^{n\times n}_\sym$.

The covariance of treelike diagrams is
defined in terms of {\em homeomorphic
matchings} between them. We start by defining this new concept.

\begin{definition}[Core]
    Let $\tau\in\calT_1$.
    Define $\textnormal{core}(\tau)$ to be
    the rooted tree obtained from
    $\tau$ by
    \begin{enumerate}
        \item Removing all hanging cactuses.
        \item Removing all non-root
        degree-2 vertices and the two edges they are incident with, and adding back a new edge between their two neighbors.
    \end{enumerate}
\end{definition}
\noindent
Note that the vertex set $V(\core(\tau))$ may be identified with a subset of $V(\tau)$, even though the second rule may lead to edges being present in $\core(\tau)$ that do not exist in $\tau$.

\begin{definition}[Homeomorphic matchings]\label{def:homeomorphic}
    Let $\tau_1,\tau_2\in\calT_1$. We
    say that a partial matching
    $P\subseteq V(\tau_1)\times V(\tau_2)$ of $\tau_1$ and $\tau_2$ is
    {\em homeomorphic} if
    \begin{enumerate}
        \item $(\textnormal{root}(\tau_1),\textnormal{root}(\tau_2))\in P$.
        \item Restricted to $V(\textnormal{core}(\tau_1))\times V(\textnormal{core}(\tau_2))$, $P$ is a rooted graph isomorphism between $\textnormal{core}(\tau_1)$ and $\textnormal{core}(\tau_2)$.

        \item Let $\{u,u'\}\in E(\textnormal{core}(\tau_1))$,
        let $(u=u_1, \ldots, u_k=u')$ be the path between $u$ and $u'$ in $\tau_1$. Let $v=P(u)$,  $v'=P(u')$, and $(v=v_1,\ldots,v_\ell=v')$ be the
        path between $v$ and $v'$ in $\tau_2$.
        Then there is no matching edge between
        $\{u_1, \ldots, u_k,v_1,\ldots,v_\ell\}$
        and its complement.
        Moreover,
        for all $(u_i,v_j)\in P$ and $(u_{i'},v_{j'})\in P$, we have $i\le i' \iff j\le j'$ (the matching restricted to the vertices in the paths is non-crossing).

        \item No inner vertices from the 
        hanging cactuses are matched.
    \end{enumerate}
    We denote by $H(\tau_1,\tau_2)$ the set
    of homeomorphic matchings between
    $\tau_1$ and $\tau_2$. 
\end{definition}

\noindent This definition
is motivated by the following lemma 
stating that, when 
computing the covariance of two treelike
diagrams, the matchings giving rise to
cactuses are precisely the homeomorphic ones.

\begin{lemma}
    \label{lem:homeomorphic}
    Let $\tau_1,\tau_2\in\calT_1$ and $\tau = \tau_1 \sqcup \tau_2$.
    For any matching $P\subseteq V(\tau_1)\times V(\tau_2)$ such that $(\roott(\tau_1), \roott(\tau_2)) \in P$, we have $\tau_P\in\calC_1$ if and only
    if $P\in H(\tau_1,\tau_2)$.
\end{lemma}

\noindent In particular, if $\tau_1, \tau_2 \in \cC_1$, only the matching $P = \{(\roott(\tau_1), \roott(\tau_2))\}$ creates a cactus $\tau_P$.
We are now ready to describe the algebra of
treelike diagrams:

\begin{lemma}\label{lem:product-gaussian}
    For all $\gamma_1,\ldots,\gamma_\ell\in\calG_1\cup \calC_1$,
    \begin{align}
        \prod_{j=1}^\ell \bz_{\gamma_j} - \sum_{M\in\calM(\ell)} \sum_{\substack{P_{uv}\in H(\gamma_u,\gamma_v)\\\forall uv\in M}} \bz_{\bigoplus_{uv\in M}\gamma_{P_{u,v}} \,\oplus\, \bigoplus_{u\notin M} \gamma_u} \in \spn(\bz_{\calA_1\setminus \calT_1})\,,\label{eq:product-gaussian}
    \end{align}
    where $\oplus$ denotes the grafting at the root.
\end{lemma}

\noindent The proofs of~\Cref{lem:homeomorphic,lem:product-gaussian} are deferred to~\Cref{sec:combinatorics}.
Note that the error in~\Cref{lem:product-gaussian}
is measured in terms of
non-treelike diagrams. 

By inverting \cref{eq:product-gaussian},
we can formulate the algebra of treelike diagrams in the language of {\em Wick products} (\cref{def:wick}).

\begin{corollary}\label{lem:treelike-hermite}
    For all $\tau\in\calT_1$,
    \begin{align}
        \bz_{\tau} - \He_{\type(\tau)}(\bz_{\calG_1}\,;\bm \Sigma) \prod_{\sigma\in\calC_1} (Z^\infty_{\sigma})^{\type(\tau)_\sigma} \in \spn(\bz_{\calA_1\setminus \calT_1})\,, \label{eq:reformulated-hermite}
    \end{align}
    where for all $\gamma,\gamma'\in\cG_1$, we defined
    the ``finite-$n$'' covariance matrix
    \begin{align}
        \bm \Sigma[\gamma,\gamma'] \defeq \sum_{P\in H(\gamma, \gamma')} \bz_{\gamma_P}\,.\label{eq:covariance-finite}
    \end{align}
\end{corollary}

\begin{proof}
    We proceed by induction on the number
    of vertices of $\tau$. First,~\cref{eq:reformulated-hermite} trivially holds
    if $\tau$ has one vertex, which proves
    the base case.
    Now, suppose that $\tau = \gamma_1\oplus \ldots \oplus \gamma_\ell$ is the
    grafting at the root of $\gamma_1,\ldots,\gamma_{\ell}\in\calG_1\cup \calC_1$.
    By~\Cref{lem:product-gaussian},
    \begin{align*}
        \bmz_\tau + \sum_{\substack{M\in \calM(\ell)\\M\neq \varnothing}} \sum_{\substack{P_{uv}\in H(\gamma_u,\gamma_v)\\\forall uv\in M}} \bz_{\bigoplus_{uv\in M}\gamma_{P_{u,v}} \,\oplus\, \bigoplus_{u\notin M} \gamma_u} - \prod_{j=1}^\ell \bmz_{\gamma_j}\in \spn(\bmz_{\calA_1\setminus \calT_1})\,.
    \end{align*}
    Applying the induction hypothesis and
    using additivity of types, we have:
    \begin{align*}
        \bz_{\bigoplus_{uv\in M}\gamma_{P_{u,v}} \,\oplus\, \bigoplus_{u\notin M} \gamma_u} - \prod_{uv\in M} \bmz_{\gamma_{P_{uv}}} \prod_{\substack{u\notin M\\\gamma_u\in \calC_1}} \bmz_{\gamma_u} \cdot \He_{\sum_{\substack{u\notin M\\\gamma_u\in\calG_1}} \type(\gamma_u)}(\bmz_{\cG_1}\,;\,\bm \Sigma)\in \spn(\bmz_{\calA_1\setminus \calT_1})\,.
    \end{align*}
    Since cactuses are not matched
    by homeomorphic matchings by
    definition, the
    product over cactuses $\bmz_{\gamma_u}$ is over all
    $u$ such that $\gamma_u\in \calC_1$, which is independent of $M$ and can be 
    factorized out. Therefore,
    in the rest of the proof
    we assume that $\gamma_i\in \calG_1$ for all $i\in [\ell]$.
    Using~\Cref{claim:non-treelike},
    we obtain
    \begin{equation}
         z_\tau + \sum_{\substack{M\in \calM(\ell)\\M\neq \varnothing}} \prod_{uv\in M} \underbrace{\sum_{P\in H(\gamma_u,\gamma_v)} \bmz_{\gamma_P}}_{\bm \Sigma[\gamma_u,\gamma_v]} \cdot  \He_{\sum_{u\notin M} \type(\gamma_u)} (\bmz_{\cG_1}\,;\, \bm \Sigma) - \prod_{j=1}^\ell \bmz_{\gamma_j}\in \spn(\bmz_{\calA_1\setminus \calT_1})\,.\label{eq:summary-wick}
    \end{equation}
    By the recursive formula of the
    Wick products (\Cref{lem:wick-recursion}),
    \begin{equation}
        \sum_{\substack{M\in \calM(\ell)\\M\neq \varnothing}} \prod_{uv\in M} \bm \Sigma[\gamma_u,\gamma_v] \cdot  \He_{\sum_{u\notin M} \type(\gamma_j)} (\bmz_{\cG_1}\,;\, \bm \Sigma) + \He_{\type(\tau)}(\bmz_{\cG_1}\,;\,\bm \Sigma)  = 
         \prod_{j=1}^\ell \bmz_{\gamma_j}\,.\label{eq:wick-application}
    \end{equation}
    Combining~\cref{eq:summary-wick,eq:wick-application} concludes the proof.
\end{proof}

Finally, if we reduce~\Cref{lem:product-gaussian} 
modulo the larger class of non-cactus diagrams
(which are the negligible diagrams {\em in expectation} under
the strong cactus property),
we deduce that the 
joint moments of the diagrams in
$\cG_1$ have an asymptotically Gaussian structure.

\begin{corollary}\label{lem:treelike-gaussian}
    For all $\gamma_1,\ldots,\gamma_\ell\in \calG_1$ and $\sigma_1,\ldots,\sigma_k\in\calC_1$,
     \[
        \prod_{i=1}^k \bmz_{\sigma_i} \left[\prod_{j=1}^\ell \bz_{\gamma_j}  - \sum_{M\in \PM(\ell)} \prod_{xy\in M} \bm \Sigma[\gamma_x,\gamma_y] \right]\in\spn(\bz_{\calA_1\setminus \calC_1})\,,
     \]
     where $\bm \Sigma$ is defined in~\cref{eq:covariance-finite}.
\end{corollary}

\begin{proof}
    Every non-treelike term in~\cref{eq:product-gaussian}
    is a fortiori not a cactus. Also, the only
    cactuses in the subtracted term occur
    when $M$ is a perfect matching. In other words,
    \[
        \prod_{i=1}^k \bmz_{\sigma_i}\left[\prod_{j=1}^\ell \bmz_{\gamma_j} - \sum_{M\in \PM(\ell)} \sum_{\substack{P_{uv}\in H(\gamma_u, \gamma_v)\\\forall uv\in M}} \bmz_{\bigoplus_{uv\in M} \gamma_{P_{uv}}}\right]\in \spn(\bmz_{\calA_1\setminus \calC_1})\,.
    \]
    Therefore, by~\Cref{claim:only-trees}, we deduce
    \[
        \bmz_{\bigoplus_{uv\in M} \gamma_{P_{uv}}} - \prod_{uv\in M} \bmz_{\gamma_{P_{uv}}}\in \spn(\bmz_{\calA_1\setminus \calC_1})\,,
    \]
    and the desired statement follows.
\end{proof}



\subsubsection{Proof of~\Cref{thm:convergence-distribution}}
\label{sec:proof-convergence}

\begin{claim}
    \label{claim:existence-joint-moments}
    Suppose that the traffic distribution
    of $\bA$ exists. Then, for any
    $\alpha_1, \ldots, \alpha_k\in\calA_1$,
    the sequence $\E \samp(\bz_{\alpha_1} \cdots \bz_{\alpha_k})$
    converges as $n\to\infty$.
\end{claim}

\begin{proof}
    This is straightforward, as
    \[
        \E \samp(\bz_{\alpha_1} \cdots \bz_{\alpha_k}) = \frac 1n \EE \sum_{i=1}^n \bz_{\alpha_1}[i] \cdots \bz_{\alpha_k}[i],
    \]
    and the inner polynomial is a scalar polynomial of $\bA$ that can be expanded in the $z$-basis of scalar diagrams as a linear combination of various quotients of the scalar diagram formed by forgetting the identity of the root in $\alpha_1 \oplus \cdots \oplus \alpha_k$.
\end{proof}

\noindent \Cref{claim:existence-joint-moments}
implies in particular 
that the sequence $\samp(\bz_{\calA_1})$ is
tight. In the rest of the proof, we 
show that the limit in distribution
actually exists and characterize it.
The following lemma is a direct consequence
of the fundamental theorem of graph
polynomials.

\begin{lemma}
    \label{lem:2ec-bounded}
    If $\|\bA\|\le O(1)$, then for each $\alpha\in\calE_1$, there
    exists $C_\alpha>0$ such that
    $|\samp(\bz_\alpha)|\le C_\alpha$.
\end{lemma}

\begin{proof}
    By~\Cref{claim:mobius} and~\Cref{lem:contract-2ec}, we can
    expand for some coefficients $c_\beta = c_\beta(\alpha)\in\R$,
    \[
        \bz_\alpha = \sum_{\beta\in\calE_1} c_\beta \bw_\beta\,.
    \]
    By~\Cref{thm:fundamental-theorem},
    it holds for every $\beta\in\calE_1$
    that $\|\bw_\beta\|_\infty\le \|\bA\|^{|E(\beta)|}$, which is at most
    $O_\alpha(1)$ by assumption.
    The lemma follows by the triangle inequality.
\end{proof}

\begin{lemma}
    \label{lem:cactus-converge}
    Suppose that the traffic distribution
    of $\bA$ exists and that~\cref{eq:tightnessNorm} holds. Then,
    $\samp(\bz_{\calC_1})$
    converges in distribution to some
    stochastic process $Z^\infty_{\calC_1}$.
\end{lemma}

\begin{proof}    
    First, assume that $\sup_{n\ge 1} \|\bm A^{(n)}\|\le K$ holds almost surely, for some universal constant $K>0$. All the moments of $\samp(\bz_{\calC_1})$ converge
    by~\Cref{claim:existence-joint-moments}.
    Since cactuses are 2-edge-connected, by~\Cref{lem:2ec-bounded}, all random
    variables $\samp(\bz_\alpha)$
    for $\alpha\in\calC_1$ are uniformly
    bounded in $n$. Hence, the moments satisfy the growth condition~\cref{eq:subexp-growth-condition}, so that
    $\samp(\bz_{\calC_1})$ converges
    in distribution by~\Cref{cor:bounded-limit-existence}. 
    Finally, if we assume~\cref{eq:tightnessNorm} rather
    than uniform boundedness, the result
    can be deduced from the latter case
    using~\Cref{lem:truncation}.
\end{proof}




\begin{proof}[Proof of~\Cref{thm:convergence-distribution}]
In the rest of the proof, we assume that
the assumptions of~\Cref{thm:convergence-distribution} are satisfied. 
We start by analyzing convergence of the subtracted term from~\Cref{lem:treelike-gaussian}. By convergence in distribution of the cactuses (\Cref{lem:cactus-converge}) and the continuous mapping theorem, we have for any 
$\gamma_1,\ldots,\gamma_\ell\in\calG_1$ and $\sigma_1,\ldots,\sigma_k\in \calC_1$,
\begin{align}
    \samp\left(\prod_{i=1}^k \bmz_{\sigma_i} \sum_{M\in \PM(\ell)} \prod_{xy\in M} \bm \Sigma[\gamma_x,\gamma_y]\right) \tod \prod_{i=1}^k Z^\infty_{\sigma_i} \sum_{M\in \PM(k)} \prod_{xy\in M} \bm \Sigma^\infty[\gamma_x,\gamma_y]\,,\label{eq:gaussian-cmt}
\end{align}
where we defined, for any $\gamma_1,\gamma_2\in\calG_1$, the ``limiting'' covariance matrix
\begin{align}
    \bm \Sigma^\infty[\gamma_1,\gamma_2] \defeq  \sum_{P\in H(\gamma_1, \gamma_2)} Z^\infty_{\gamma_P}\,.\label{eq:covariance}
\end{align}
Since all joint moments converge
by~\Cref{claim:existence-joint-moments},
the sequence
of random
variables on the left-hand side of~\cref{eq:gaussian-cmt} is uniformly
integrable. So we also get convergence 
of the mean,
\[
    \E \samp\left(\prod_{i=1}^k \bmz_{\sigma_i} \sum_{M\in \PM(\ell)} \prod_{xy\in M} \bm \Sigma[\gamma_x,\gamma_y]\right) \underset{n\to\infty}\longrightarrow \E\left[\prod_{i=1}^k Z^\infty_{\sigma_i} \sum_{M\in \PM(k)} \prod_{xy\in M} \bm \Sigma^\infty[\gamma_x,\gamma_y]\right]\,.
\]   
Combining with~\Cref{lem:treelike-gaussian} and the strong cactus property,
\begin{align}
    \E \samp\left(\prod_{i=1}^k \bmz_{\sigma_i} \prod_{j=1}^\ell \bz_{\gamma_j}\right) \underset{n\to\infty}\longrightarrow \E\left[ \prod_{i=1}^k Z^\infty_{\sigma_i} \sum_{M\in \PM(\ell)} \prod_{xy\in M} \bm \Sigma^\infty[\gamma_x,\gamma_y]\right]\,.\label{eq:covariance-converge}
\end{align}
The right-hand side of~\cref{{eq:covariance-converge}}
coincides with the moments of $Z^\infty_{\calG_1\cup \calC_1}$. Recall that the law of
$Z^\infty_{\calG_1}$ satisfies that after sampling
$Z^\infty_{\calC_1}$ from its marginal
(which is bounded almost surely by~\Cref{lem:2ec-bounded}), then $Z^\infty_{\calG_1}$
conditioned on $Z^\infty_{\calC_1}$ is a Gaussian
process with covariance kernel given by~\cref{eq:covariance}. This
object satisfies the moment growth
condition~\cref{eq:subexp-growth-condition}. So~\Cref{cor:bounded-limit-existence}
applies and we obtain convergence in
distribution of $\samp(\bz_{\calG_1\cup \calC_1})$
to $Z^\infty_{\calG_1\cup \calC_1}$.


By~\Cref{lem:non-treelike}, the non-treelike diagrams converge in $L^2$ to 0,
so by Slutsky's lemma, we obtain joint
convergence in distribution, except for the remaining treelike, non-Gaussian trees. By~\Cref{lem:treelike-hermite}, these are continuous images
of cactuses and non-treelike diagrams, so
by the continuous mapping theorem, all
diagrams
converge jointly in distribution to $Z^\infty_{\calA_1}$.
\end{proof}



\subsection{The treelike AMP algorithm}

Now we turn to studying the dynamics of GFOM operations.

\begin{definition}[Asymptotic state]\label{def:asymptotic-state}
    Let $(\bm x_i)_{i\in \calI}$ be a
    family of random vectors,
    $\bm x_i\in \R^{n}$.
    We say that a stochastic process $(X_i)_{i\in \calI}$
    is the asymptotic state of $(\bm x_i)_{i\in \calI}$ if, for any $k\ge 1$, $i_1, \ldots, i_k\in \calI$, and any bounded continuous
    or polynomial function $\varphi:\R^k\to \R$,
    \begin{equation}
        \lim_{n\to\infty} \frac 1n \sum_{j=1}^n \E\varphi(\bm x_{i_1}[j], \ldots, \bm x_{i_k}[j]) = \E\varphi(X_{i_1}, \ldots, X_{i_k})\,.\label{eq:defAsymptoticState}
    \end{equation}
\end{definition}

\noindent \Cref{def:asymptotic-state} requires in particular for $(\bm x_i)_{i\in \mathcal I}$ to converge in distribution to $(X_i)_{i\in \calI}$. 
As with convergence in distribution in general, this suffers from the caveat that the law of the limit in distribution of $(X_i)_{i\in \calI}$ is unique, but the probability space on which the limit $(X_i)_{i\in \calI}$ is realized is not.
Thus when we speak of ``the asymptotic state'' we refer to a specific law, not a specific collection of random variables.
Nonetheless, the sampling procedure in \Cref{thm:convergence-distribution} suggests a natural
way to sample an asymptotic state of the
iterates of a pGFOM, since, provided we know how to sample from $Z_{\calC_1}^{\infty}$ (which we must address on a case-by-case basis), the other $Z_{\alpha}^{\infty}$ are conditionally Gaussian or deterministic functions thereof.

Translating the limiting variables $Z_{\alpha}^{\infty}$ from \Cref{thm:convergence-distribution} to a construction of an asymptotic state, we find:
\begin{lemma}\label{claim:cvImpliesState}
    Assume that $\bA = \bA^{(n)}$ satisfies the assumptions of~\Cref{thm:convergence-distribution}.
    Let 
    \begin{equation}
        \bm x = \sum_{\alpha\in \calA_1} c_\alpha \bmz_\alpha(\bmA)\label{eq:assumeState}
    \end{equation}
    for
    some finitely supported coefficients
    $(c_\alpha)_{\alpha\in \calA_1}$. Then,
    \begin{equation}
        X \defeq \sum_{\alpha\in \calA_1} c_\alpha Z^\infty_\alpha\label{eq:asymptState}
    \end{equation}
    is the asymptotic state of $\bm x$.
    Moreover, if $\bm x_t$ is of the form~\cref{eq:assumeState} for any $t\ge 1$ and $X_t$
    is correspondingly defined as in~\cref{eq:asymptState}, then $(X_t)_{t\ge 1}$ is the asymptotic
    state of $(\bm x_t)_{t\ge 1}$. 
\end{lemma}

\noindent We emphasize that the index set
$t\ge 1$ is independent of $n$, and so our
results hold for all fixed iterates $t$
independent of $n$, in the limit $n\to\infty$.

\begin{proof}
    The statement for bounded continuous test
    functions $\varphi$ follows from~\Cref{thm:convergence-distribution} and
    the continuous mapping theorem.
    For polynomial $\varphi$, we proceed by a truncation argument.
    Let $S_n \defeq \samp(\bm x_1, \ldots, \bm x_t)$ and $S \defeq (X_1, \ldots, X_t)$. Fix a cutoff $K>0$ and consider
    any bounded continuous function $\varphi_K$ such
    that $|\varphi_K|\le |\varphi|$, $\varphi_K(s) = \varphi(s)$ for all
    $\|s\|_2\le K$ and $\varphi_K(s)=0$ for all
    $\|s\|_2>2K$ (standard approximations show that such a function exists). First, $\left|\E \varphi_K(S_n)-\E \varphi_K(S)\right|$ converges to $0$
    as $n\to\infty$
    by the bounded continuous case. Next,
    \begin{align*}
        \left|\E \varphi(S_n) - \E \varphi_K(S_n)\right|&\le \E \left[|\varphi(S_n)| \mathbf 1_{\|S_n\|_2>K}\right]\tag*{(Definition of the truncated function)}\\
        &\le (\E \varphi(S_n)^2)^{\frac 12} \Pr(\|S_n\|_2>K)^{\frac 12}\tag*{(Cauchy-Schwarz inequality)}\\
        &\le (\E \varphi(S_n)^2)^{\frac 12} \frac {(\E \|S_n\|_2^2)^{\frac 12}} {K}\tag*{(Markov inequality)}
    \end{align*}
    Note that these quantities are respectively equal to
    \[
        \E \varphi(S_n)^2 = \frac 1n \sum_{i=1}^n \varphi(\bm x_1[i], \ldots, \bm x_t[i])^2\quad\text{and} \quad\E \|S_n\|_2^2 = \sum_{s=1}^t \frac 1n  \sum_{i=1}^n \bm x_s[i]^2\,,
    \]
    which both converge as $n\to\infty$ by existence
    of the traffic distribution, and in particular are bounded uniformly in $n$. Hence, there exists $C>0$ independent of $n$ and $K$ such that
    \[\limsup_{n\to\infty} \left|\E \varphi(S_n) - \E \varphi_K(S_n)\right|\le \frac C K\,,\] and the
    same holds for $\E \varphi(S)-\E \varphi_K(S)$ 
    by the same argument, using the fact that all moments exist in the space generated by $Z_{\calA_1}^\infty$. We obtain $\limsup_{n\to\infty} \left|\E \varphi(S_n) - \E \varphi_K(S_n)\right|\le 2C/K$, and the claim
    follows from taking the limit $K\to\infty$.
\end{proof}



By~\Cref{prop:gfom-polynomial}, the iterates
of a pGFOM are of the form~\cref{eq:assumeState}, so they have an asymptotic state. By definition, these asymptotic states determine the
limiting distribution of any (bounded continuous or polynomial) observable.
Motivated by this,
we introduce a family of approximate message passing algorithms whose
asymptotic states are conditionally Gaussian.

\begin{theorem}[Treelike AMP]
\label{thm:full-onsager}
Assume that $\bA=\bA^{(n)}$ satisfies the assumptions of \Cref{thm:convergence-distribution}. 
Let $f_t : \R \to \R$ be polynomial functions.\footnote{For ease of exposition $f_t$ is assumed to be ``memoryless'', meaning that it only takes the most recent $\bx_{t}$ as input.} Define:
\begin{equation}
    \begin{aligned}
    \bmx_0 &= \bm{1}\,,\qquad 
    \bmx_t = \bA \bmf_{t-1}
      - \sum_{s=0}^{t-1} \bmb_{s,t} \cdot \bmf_s\,,\\
    \bmb_{s,t}[i]
    &\defeq \sum_{\substack{i_s,\ldots,i_{t-1}\in[n] \textnormal{ distinct}\\i_s=i}} \left(\prod_{r=s+1}^{t-1} \bm A[i_{r-1}, i_{r}] \bm f'_{r}[i_{r}] \right)\bA[i_{t-1},i_s]\,,\\
    \bmf_t &\defeq f_t(\bmx_t)\,, \qquad
    \bmf'_t \defeq f'_t(\bmx_t)\,,\qquad \bm f_0 = \bm 1\,.\label{eq:tree-amp}
    \end{aligned}
\end{equation}
Then $\bx_t \in \linspan(\bmz_{\cG_1 \cup (\cA_1 \setminus \cT_1)}(\bA))$. Therefore, the asymptotic state of $(\bm x_t)_{t\ge 1}$ defined
in~\cref{eq:asymptState}
is a centered Gaussian process conditionally on $Z^\infty_{\calC_1}$.
\end{theorem}

\noindent To prove~\Cref{thm:full-onsager}, motivated by the results in \cref{sec:vector-limit}, we introduce
the following handy notations:
\newcommand\gauss{\textnormal{gaussian}}

\begin{definition}[Equality modulo non-treelike diagrams]\label{def:eqinf}
    For $\bm x,\bm y\in \spn(\bmz_{\calA_1})$, we write $\bm x \eqinf \bm y$ if $\bm x - \bm y\in \spn(\bmz_{\calA_1\setminus \calT_1})$.
    We denote by $\cactus(\bmx)$ the
    projection of $\bmx$ onto the span
    of the cactus diagrams $\cC_1$,
    and by $\gauss(\bmx)$ the
    projection of $\bmx$ onto the span
    of the Gaussian diagrams $\cG_1$.
\end{definition}


The iterates of the treelike AMP
algorithm~\cref{eq:tree-amp} are engineered to 
asymptotically generate a self-avoiding walk.
That is, whenever the algorithm performs a matrix multiplication operation,
the Onsager correction terms in \cref{eq:tree-amp} (the subtracted terms involving $\bb_{s,t}$) are chosen to subtract
off the terms in the resulting diagram expansion which (1)~are treelike and (2)~revisit an existing vertex in any diagram.

\begin{example}[Self-avoiding walk]
For intuition, consider
the case of~\Cref{thm:full-onsager}
where $f_t(x) = x$.
Let $\pi_t$ be the $t$-path
diagram and $\rho_t$ the $t$-cycle
diagram. We can expand exactly:
\begin{align*}
    \bA \bmz_{\pi_t} &= \bmz_{\pi_{t+1}} + \sum_{s = 0}^t \bmz_{\rho_{s+1} \oplus \pi_{t-s}}\,.
\end{align*}
For each term on the right-hand side, we have the approximate factorization (by~\Cref{lem:product-gaussian}) $\bmz_{\rho_{s+1} \oplus \pi_{t-s}} \eqinf \bmz_{\rho_{s+1}} \cdot \bmz_{\pi_{t-s}}$,
which holds up to non-treelike terms.
Then, we define a self-avoiding version of power iteration by:
\[
    \bmx_0 = \bm{1}, \qquad \bmx_{t+1} = \bA \bmx_t - \sum_{s = 0}^t \bmz_{\rho_{s+1}} \cdot \bmx_{t-s}\,.
\]
By construction, we have $\bmx_t \eqinf \bmz_{\pi_t}$ and therefore, assuming the conditions of \cref{thm:convergence-distribution}, the asymptotic state $X_t$ of $\bm x_t$ is Gaussian.
\end{example}

To analyze a general iteration in the proof of~\Cref{thm:full-onsager}, we separate the diagram expansion of $\bmf_t$ into its linear and nonlinear parts:
\begin{equation*}\label{eq:linear-nonlinear}
    \bmf_t = \underbrace{\sum_{\gam \in \cG_1} c_\gam \bmz_\gam(\bA)}_{=: \bmf^1_t} + \underbrace{\sum_{\tau \in \cT_1 \setminus \cG_1} c_\tau \bmz_\tau(\bA)}_{=: \bmf^{\neq 1}_t} + \sum_{\al \in \cA_1 \setminus \cT_1} c_\al \bmz_\al(\bA)\,.
\end{equation*}

\noindent We call $\bmf^1_t = \gauss(\bmf_t)$ the ``linear part'' since it should be thought of as the degree-1 part of the Hermite expansion
of $\bmf_t$ with respect to the Gaussian vectors $\bmz_{\cG_1}(\bA)$,
while $\bmf^{\neq 1}_t$ equals all of the other components of the Hermite expansion.
More precisely, when $\bmf_t$ is of the form $f_t(\bmx_{t})$ for some Gaussian vector $\bmx_{t}$, which is the situation for AMP, then $\bmf^{1}_t$ has the following simple form.

\begin{lemma}\label{lem:f-linear}
    Let $\bmx \in \linspan(\bmz_{\cG_1})$ and let $f : \R \to \R$ be a polynomial.
    Then,
    \[
        \gauss(f(\bm x)) \eqinf \cactus(f'(\bm x)) \cdot \bm x\,.
    \]
\end{lemma}
\begin{proof}
    Suppose that $f(x) = x^\ell$ for some integer $\ell\ge 0$ (the general case follows by linearity). By~\Cref{lem:product-gaussian},
    the product of
    diagrams in $\calG_1$
    yields a diagram in $\calG_1$ only
    when every diagram except one is matched. Formally, write $\bmx = \sum_{\gamma\in\calG_1} c_\gamma \bmz_{\gamma}$, then by~\Cref{lem:product-gaussian},
    \begin{align}
        \bmf^1 \eqinf \ell \sum_{\gamma_1,\ldots,\gamma_\ell\in \calG_1} c_{\gamma_1}\ldots c_{\gamma_{\ell}} \sum_{M\in \PM(\ell-1)} \sum_{\substack{P_{uv}\in H(\gamma_u,\gamma_v)\\\forall uv\in M}} \bz_{\bigoplus_{uv\in M}\gamma_{P_{u,v}} \oplus \gamma_\ell} \,.\label{eq:l-power}
    \end{align}
    Viewing $\bigoplus_{uv\in M}\gamma_{P_{u,v}}$ as a fixed cactus, by~\Cref{lem:product-gaussian},
    every term on the right-hand side
    satisfies
    \begin{align}
        \bz_{\bigoplus_{uv\in M}\gamma_{P_{u,v}} \oplus \gamma_\ell} \eqinf \bz_{\bigoplus_{uv\in M}\gamma_{P_{u,v}}} \cdot \bmz_{\gamma_\ell}\,.\label{eq:factorize-inter}
    \end{align}
    Applying~\Cref{lem:product-gaussian}
    one last time, it remains to observe 
    that the cactus part of $f'(\bm x) = \ell \bm x^{\ell-1}$ is
    \begin{align}
        \cactus(f'(\bmx)) \eqinf \ell \sum_{\gamma_1,\ldots,\gamma_{\ell-1}\in \calG_1} c_{\gamma_1} \ldots c_{\gamma_{\ell-1}} \sum_{M\in \PM(\ell-1)} \sum_{\substack{P_{uv}\in H(\gamma_u,\gamma_v)\\\forall uv\in M}} \bz_{\bigoplus_{uv\in M}\gamma_{P_{u,v}}}\,.\label{eq:l-power-minus}
    \end{align}
    Combining~\cref{eq:l-power,eq:factorize-inter,eq:l-power-minus} yields the desired claim.
\end{proof}

The next key~\cref{lem:memory-formula} derives an explicit asymptotic formula for the AMP iterates: $\bmx_t$ is generated by taking a self-avoiding walk from each nonlinear term $\bmf_s^{\neq 1}$.

\begin{definition}\label{def:self-avoiding}
Let $\bmc_t = \cactus(f'_t(\bmx_t))$.
Define the \emph{self-avoiding walk matrix} $\bmB_{s,t}$ generated by the iteration between time $s$ and $t$ to be:
\begin{align*}
    \bmB_{s,t}[i,j] &\defeq \sum_{\substack{i_{s},\dots, i_{t} \in [n] \textnormal{ distinct}\\i_{s} = j,\; i_{t} = i}} 
    \bmA[i_s, i_{s+1}] \cdots \bmA[i_{t-1},i_t] \cdot \bm c_{s+1}[i_{s+1}] \cdots \bm c_{t-1}[i_{t-1}]\,.
\end{align*}
\end{definition}

\noindent Recalling \Cref{def:open-cactus}, $\bmB_{s,t}$
is a linear combination of
open cactus matrices in the $z$-basis (up to non-treelike terms which arise from intersections involving the $\bm c_{s+i}$).
For example, $\bmB_{t-1,t}$ equals $\bA$ with the diagonal elements set to zero.
We note the analogy between $\bmb_{s,t}$ and $\bmB_{s,t}$, which contain similar self-avoiding walks that return to the start and do not return to the start, respectively.

\begin{lemma}\label{lem:memory-formula}
    Define $\bmx_t, \bmf_t$ by \cref{eq:tree-amp} and let $\bmc_t = \cactus(f'_t(\bmx_t))$. Then for $t \geq 1$:
    \begin{align*}
        \bmx_t \eqinf \sum_{s = 0}^{t-1} \bmB_{s,t} \bmf_s^{\neq 1} \quad \text{and} \quad \bmf_t \eqinf \bmc_t \cdot \sum_{s=0}^{t-1}  \bmB_{s,t} \bmf_s^{\ne 1} + \bmf_{t}^{\neq 1}\,.
    \end{align*}
\end{lemma}
\begin{proof}
    First, note that for a fixed $t$, the second equation follows from the first:
    \begin{align*}
        \bmf_t &\eqinf \bmf_t^1 + \bmf_t^{\neq 1}\\
        &\eqinf \bc_t  \cdot \bmx_t + \bmf_t^{\neq 1}  \tag{\cref{lem:f-linear}} \\
        &\eqinf \bc_{t} \cdot \sum_{s = 0}^{t-1}  \bmB_{s,t} \bmf_s^{\neq 1} + \bmf_{t}^{\neq 1} \tag{first equation and~\Cref{claim:only-trees}}
    \end{align*}
    To establish the equations, we use induction on $t$. In the base case $t = 1$ we have $\bmx_1 = \bA \bmf_0 - \bmb_{0,1}\cdot \bmf_0 = \bmB_{0,1} \bmf_0$ as needed.
    Now, assume that the formulas
    hold for $0, \ldots, t$.
    Denote by $\bmC_t$ the diagonal matrix with entries $\bmc_t$.   
    The equation for $\bmf_t$ implies:
    \begin{align*}
        \bA \bmf_{t} &\eqinf \sum_{s = 0}^{t-1} \bA \bmC_{t} \bmB_{s, t} \bmf_s^{\neq 1}  + \bA \bmf_{t}^{\neq 1} \numberthis \label{eq:Af_t}
    \end{align*}
    If we expand the matrix product $\bA \bC_{t} \bB_{s,t}$ we can partition the sum based on whether the matrix $\bA$ revisits a vertex already on the walk:
    \begingroup
    \begin{align}
        (\bA\bC_{t} \bB_{s, t})[i,j] &= \sum_{k=1}^n \bmA[i,k] \bm c_t[k] \sum_{\substack{i_{s},\dots, i_{t} \in [n] \textnormal{ distinct}\\i_{s} = j,\; i_{t} = k}} 
    \bmA[i_s, i_{s+1}] \cdots \bmA[i_{t-1},i_t] \cdot \bm c_{s+1}[i_{s+1}] \cdots \bm c_{t-1}[i_{t-1}]\nonumber\\
        &= \sum_{\substack{i_{s},\dots, i_{t+1} \in [n] \textnormal{ distinct}\\i_{s} = j,\; i_{t+1} = i}} 
    \bmA[i_s, i_{s+1}] \cdots \bmA[i_{t},i_{t+1}] \cdot \bm c_{s+1}[i_{s+1}] \cdots \bm c_{t}[i_{t}]\label{eq:first-term-gaussian}\\
    &+ \sum_{r=s}^t \sum_{\substack{i_{s},\dots, i_{t} \in [n] \textnormal{ distinct}\\i_{s} = j,\; i_{r} = i}} 
    \bmA[i_s, i_{s+1}] \cdots \bmA[i_{t-1},i_t] \cdot \bmA[i_r,i_t] \cdot \bm c_{s+1}[i_{s+1}] \cdots \bm c_{t}[i_{t}]\,.\label{eq:second-term-gaussian}
    \end{align}

    The first term~\cref{eq:first-term-gaussian} is self-avoiding and equals $\bmB_{s,t+1}[i,j]$.
    In the second term~\cref{eq:second-term-gaussian}, the term $r=s$ is diagrammatically a cycle
    and is equal to $\bmb_{s,t+1}[i]$ when $i=j$, and $0$ otherwise:

    \begin{claim}\label{lem:bst-c}
    We have:
    \[
        \bmb_{s,t} \eqinf \left(\sum_{\substack{i_s,\ldots,i_{t-1}\in[n] \textnormal{ distinct}\\i_s=i}}
        \bA[i_s,i_{s+1}] \cdots \bA[i_{t-2},i_{t-1}] \bA[i_{t-1},i_s] \cdot \bmc_{s+1}[i_{s+1}] \cdots \bmc_{t-1}[i_{t-1}]\right)_{i\in [n]}\,.
    \]
    \end{claim}
    \begin{proof}
        The only difference between this formula and the definition of $\bb_{s,t}$ is that the vectors $\bmf'_t$ at the internal vertices of the cycle have been replaced by $\bmc_t$.
        This holds up to non-treelike terms since placing a non-cactus diagram at any internal vertex of the cycle will create only non-treelike diagrams.
    \end{proof}    
    \noindent The remaining terms in~\cref{eq:second-term-gaussian} are a cycle and a path joined together at vertex $r$:
    
    \begin{claim}\label{claim:factorize-inside-path-cycle}
        Let $r\in \{s+1, \ldots, t\}$. For $i\in [n]$, let
        \[
            \bm u[i] =         \sum_{\substack{i_{s},\dots, i_{t} \in [n] \textnormal{ distinct}\\i_{r} = i}} 
    \bmA[i_s, i_{s+1}] \cdots \bmA[i_{t-1},i_t] \cdot \bmA[i_r,i_t] \cdot \bm c_{s+1}[i_{s+1}] \cdots \bm c_{t}[i_{t}] \bmf_s^{\neq 1}[i_s]\,.
        \]
        Then $\bm u \eqinf \bm b_{r,t+1} \cdot \bm C_r \bm B_{s,r} \bmf_s^{\neq 1}$.
    \end{claim}

    \begin{proof}
        By expanding definitions, we can conveniently interpret
        \[
            \bm b_{r,t+1} \cdot \bm C_r \bm B_{s,r} \bmf_s^{\neq 1}[i] = \sum_{\substack{i_{s},\dots, i_{t} \in [n]\\
            i_s,\ldots,i_{r} \textnormal{ distinct}\\
            i_r,\ldots,i_{t}\textnormal{ distinct}\\i_{r} = i}} 
    \bmA[i_s, i_{s+1}] \cdots \bmA[i_{t-1},i_t] \cdot \bmA[i_r,i_t] \cdot \bm c_{s+1}[i_{s+1}] \cdots \bm c_{t}[i_{t}] \bmf_s^{\neq 1}[i_s]\,.
        \]
        Since the diagram induced on
        $\{i_r, \ldots, i_t\}$ is
        a cycle, any intersection between the
        vertices $\{i_s, \ldots, i_r\}$
        and $\{i_r, \ldots, i_t\}$
        would create a non-treelike 
        diagram.
    \end{proof}
    
    Plugging~\Cref{claim:factorize-inside-path-cycle} in to~\cref{eq:Af_t}, we have:
    \begin{align*}
        \bA \bmf_t &\eqinf \sum_{s = 0}^{t-1} \bB_{s,t+1} \bmf_s^{\neq 1} + \sum_{s=0}^{t-1}\bmb_{s,t+1} \cdot \bmf_s^{\neq 1} + \sum_{s=0}^{t-1}\sum_{r=s+1}^t \bmb_{r,t+1}\cdot (\bmC_r \bmB_{s,r} \bmf^{\neq 1}_s)+ \underbrace{\bA \bmf_t^{\neq 1}}_{= \bB_{t,t+1} \bmf_t^{\neq 1} + \bmb_{t,t+1} \cdot \bmf_t^{\neq 1}}\\
        &= \sum_{s = 0}^{t} \bB_{s,t+1} \bmf^{\neq 1}_s + \sum_{r = 0}^t \bmb_{r,t+1} \cdot \left(\bmC_r \sum_{s = 0}^{r-1} \bmB_{s,r}\bmf_s^{\neq 1} + \bmf_r^{\neq 1} \right)\\
        &\eqinf \sum_{s = 0}^{t} \bB_{s,t+1} \bmf^{\neq 1}_s + \sum_{r = 0}^t \bmb_{r,t+1} \cdot \bmf_r
    \end{align*}
    \endgroup
    The last equality is the inductive formula for $\bmf_r$.
    The Onsager correction subtracts off the second sum, leaving only the desired first sum for $\bmx_{t+1}$.
\end{proof}


\begin{proof}[Proof of \cref{thm:full-onsager}.]
    We prove the following purely combinatorial claim about~\cref{eq:tree-amp}:
    $\bm x_t$ is in
    the span of non-treelike
    diagrams and {\em Gaussian} treelike
    diagrams.  
    By~\Cref{claim:cvImpliesState}, this will
    imply
    that conditioned on $Z^\infty_{\calC_1}$, the asymptotic state of $(\bm x_t)_{t\ge 1}$ is a centered Gaussian process, as desired.

    To show the claim, 
    we start from the conclusion of~\cref{lem:memory-formula}:
    \[
        \bmx_t \eqinf \sum_{s=0}^{t-1}\bB_{s,t}\bmf_s^{\neq 1}\,.
    \]
    The diagrams in $\bB_{s,t}\bmf_{s}^{\neq 1}$ are obtained by: (1)~choose a diagram from $\bmf_{s}^{\neq 1}$, (2)~choose a cactus diagram from $\bc_r$ at each internal vertex of $\bB_{s,t}$
    (i.e.\ each internal vertex along a path of length $t-s$),
    (3)~multiply these diagrams together.
    Since none of the diagrams in $\bmf_s^{\neq 1}$ have degree 1 at the root by definition,
    the only treelike terms in the product are formed by grafting the diagrams together without intersections.
    In particular, the root is the endpoint of the path in $\bB_{s,t}$ and has degree 1.
    This concludes the proof.
\end{proof}

\subsubsection{Covariance structure of treelike AMP}

While~\Cref{thm:full-onsager} shows that the treelike AMP iterates are
asymptotically Gaussian, it does not identify their covariance. We calculate the covariance ``combinatorially'' by calculating the cactus diagrams appearing in
the expansion of $\langle \bmx_s, \bmx_t\rangle$.


\begin{proposition}\label{lem:cov-cactus}
    Let $\bm x_t$ follow the iteration~\cref{eq:tree-amp}. Then for any $s,t\ge 1$,
    \begin{equation}\label{eq:cov-finite}
        \bm x_s \cdot \bm x_t - \sum_{s'=0}^{s-1} \sum_{t'=0}^{t-1} \bm B_{s'st't} (\bm f_{s'} \cdot \bm f_{t'}) \in \spn(\bm z_{\calA_1\setminus \calC_1})\,,
    \end{equation}
    where for $0\le s'\le s, 0\le t'\le t$, we define the matrix $\bmB_{s'st't} \in \R^{n\times n}$ by
    \begin{equation*}
        \bm B_{s'st't}[i,j] \defeq \sum_{\substack{i_{s'},\ldots,i_{s}, j_{t'}, \ldots, j_{t}\in [n]\\
        \textnormal{distinct except}\\
        i_{s'}=j_{t'}=j,\, i_s=j_t=i}} \prod_{r=s'}^{s-1} \bm A[i_r,i_{r+1}] \prod_{r=t'}^{t-1} \bm A[j_r,j_{r+1}] \prod_{r=s'+1}^{s-1} \bm c_r[i_r] \prod_{r=t'+1}^{t-1} \bm c_r[j_r]\,.
    \end{equation*}
\end{proposition}


When we apply this lemma, we will average \cref{eq:cov-finite}
over the coordinates $i \in [n]$ and over $\bA$.
Since the error terms are all in the span of non-cactus diagrams,
all error terms converge to 0 by the strong cactus property.
On the other hand,
the average of the term $\bmx_s \cdot \bmx_t$ converges to the covariance $\E[X_sX_t]$
which we want to calculate.
The subtracted terms involving the $\bB_{s't'st}$ matrices converge to limits depending on the asymptotic values of the cactuses $Z_{\cC_1}^\infty$.
For some settings (such as \cref{prop:factorizing-deterministic}), the values $Z_{\cC_1}^\infty$ are deterministic.
For other settings, the values $Z_{\cC_1}^\infty$ are random, and we will condition on them in order to obtain the conditional covariance.


\begin{proof}
    Since $\bm x_s$ and $\bm x_t$ have
    degree exactly one at the root, in
    order to form a cactus in $\bm x_s\cdot \bm x_t$, the paths
    from the root of $\bm x_s$ and $\bm x_t$ in the expansion from~\Cref{lem:memory-formula} must meet at some point. This
    intersection
    cannot happen at a vertex from $\bm f_{s'}^{\neq 1}$ or $\bm f_{t'}^{\neq 1}$ (that would create edges in two cycles).
    Let
    $s'\in \{0,\ldots,s-1\}$ and $t'\in \{0,\ldots,t-1\}$ denote the integers
    such that the first intersection
    corresponds to the indices
    $i_{s'}$ (for $\bm x_s$) and $i_{t'}$ (for $\bm x_t$) in~\Cref{def:self-avoiding}. Then,
    we can decompose
    \[
        \bm x_s\cdot \bm x_t - \sum_{s'=0}^{s-1} \sum_{t'=0}^{t-1} \bm B_{s'st't} ((\bm c_{s'} \cdot \bm x_{s'} + \bm f_{s'}^{\neq 1}) \cdot (\bm c_{t'} \cdot \bm x_{t'} + \bm f_{t'}^{\neq 1}))\in \spn(\bm z_{\calA_1\setminus \calC_1})\,,
    \]
    and the conclusion follows from
    the equality (\Cref{lem:f-linear}) $\bm f_s\eqinf \bm c_s \cdot \bm x_s + \bm f_s^{\neq 1}$.
\end{proof}

\noindent The 
cactus expansion of $\bm B_{s'st't} (\bm f_{s'} \cdot \bm f_{t'})$ can be obtained explicitly by
combining a cycle of length $s-s'+t-t'$
along the edges of $\bm B_{s'st't}$,
a cactus from $\bm c_r$ hanging at every vertex $r$ in the cycle, and a homeomorphic
matching of the tree components of $\bm f_{s'}$ and $\bm f_{t'}$ (\Cref{def:homeomorphic}). 

\subsection{Examples of state evolution}\label{sec:se-examples}

In this section, we specialize~\cref{thm:full-onsager}
to obtain a more explicit description
of the state evolution of the
treelike AMP algorithm for several concrete matrix models. 

\begin{notation}
    \label{not:empirical}
    For a vector $\bm x\in \R^n$, we will use the following notation for empirical averages:
    \[
        \langle \bm x \rangle \defeq \frac 1n \sum_{i=1}^n x_i\,.
    \]
\end{notation}

Technically, most algorithms in this section are not pGFOM since they
calculate empirical averages.
Assuming that the traffic distribution concentrates and the vector $\bm x$ lies in the diagram basis,
then the empirical average $\langle{\bm x}\rangle$ concentrates,
and we can replace $\langle{\bm x}\rangle$ by its limit $\E X$ without changing the asymptotic state of the algorithm.
This is formally proven in \cref{lem:asymptotic-state-empirical-average}.


\subsubsection{Orthogonally invariant random matrices}
\label{sec:treelike-oamp}

In the special case that $\bA$ is drawn from an orthogonally invariant random matrix ensemble,
the treelike AMP algorithm recovers the orthogonal AMP algorithm of Fan~\cite{fan2022approximate},
giving a new proof of this result.

\begin{theorem}[State evolution for orthogonally invariant matrices]\label{thm:orthogonal-amp}
    Let $\bA = \bA^{(n)} \in \R_\sym^{n \times n}$ be an orthogonally invariant random matrix converging in tracial moments in $L^2$ to a probability measure with free cumulants $(\kappa_q)_{q\ge 1}$. Assume $\bm A$ satisfies~\cref{eq:tightnessNorm}.
    Let $f_t : \R \to \R$ be polynomial functions and define the iteration
    \begin{alignat*}{2}
        \bmx_0 &= \bm{1}, \qquad &\bmx_t &= \bA f_{t-1}(\bm x_{t-1}) - \sum_{s = 0}^{t-1} \kappa_{t-s} \left(\prod_{r=s+1}^{t-1} \langle f'_r(\bm x_r)\rangle \right)f_s(\bm x_s)\quad \forall t\ge 1\,. \numberthis\label{eq:rotation-amp}
    \end{alignat*}
    Then,
    the asymptotic state of $(\bm x_t)_{t\ge 1}$ is a centered Gaussian process $(X_t)_{t\ge 1}$ with covariance
    \begin{align*}
        \E\left[X_s X_t\right] &= \sum_{s' = 0}^{s-1} \sum_{t' = 0}^{t-1} \kappa_{s-s'+t-t'} \left(\prod_{r = s'+1}^{s-1} \E f'_r(X_r) \right) \left(\prod_{r=t'+1}^{t-1} \E f'_r(X_r)\right) \E\left[f_{s'}(X_{s'}) f_{t'}(X_{t'})\right]\quad \forall s,t\ge 1\,,
    \end{align*}
    with $X_0\defeq 1$.
\end{theorem}

\begin{proof}
    By~\cref{thm:moments}, $\bA$ satisfies the factorizing strong cactus property and its diagonal distribution exists,
    so the assumptions of~\Cref{thm:convergence-distribution} and \Cref{thm:full-onsager} are satisfied.
    Therefore, the treelike AMP algorithm in \cref{eq:tree-amp} has Gaussian asymptotic state.
    
    We now specialize the Onsager correction term in~\cref{eq:tree-amp} to
    this model. The term $\bm b_{s,t}$ is represented by a cycle of length
    $t-s$, with $f'_r(\bm x_r)$ attached to the $r$th vertex of the cycle for each
    $s<r<t$. By~\Cref{claim:only-trees}, we
    only need to look at treelike contributions in $\bm b_{s,t}$.
    Because
    of the base cycle, these are only cactuses, obtained by attaching
    cactuses from $(f'_r(\bm x_r))_{s<r<t}$
    along the base cycle. 
    By~\Cref{prop:factorizing-deterministic},
    $\bm b_{s,t}$ has constant asymptotic
    state equal to
    $\kappa_{t-s}\prod_{r=s+1}^{t-1}\E f'_r(X_r)$.
    The cactuses in $\bm b_{s,t}$ persist until the end of the algorithm,
    so that they will eventually contribute this value towards the asymptotic state.
    Hence it does not affect the asymptotic state to replace $\bm b_{s,t}$ immediately by its limiting constant value.
    
    Moreover, by~\Cref{lem:asymptotic-state-empirical-average} and \cref{lemma:full-factorization}, we may replace $\E f'_r(X_r)$ by the empirical average
    $\langle f'_r (\bm x_r)\rangle$ to obtain~\cref{eq:rotation-amp} without affecting the asymptotic state.
    Now the asymptotic state $X_t$ of~\cref{eq:rotation-amp} matches that of~\cref{eq:tree-amp}, and we may apply~\cref{thm:full-onsager} to deduce that $X_t$ is Gaussian.

    To calculate the covariance $\E\left[X_sX_t\right]$, we average \cref{lem:cov-cactus} over the coordinates $i \in [n]$ and take the limit $n \to \infty$.
    On the right side of \cref{eq:cov-finite},
    the cycle of $\bB_{s'st't}$ contributes $\kappa_{s-s'+t-t'}$ and the hanging diagrams $f'_r(\bm x_r)$ inside $\bB_{s'st't}$ contribute
    $\E f'_r(X_r)$ by the factorizing cactus property.
    The cactuses in $f_{s'}(\bm x_{s'}) \cdot f_{t'}(\bm x_{t'})$ contribute $\E\left[f_{s'}(X_{s'})f_{t'}(X_{t'})\right]$, which establishes the desired recurrence.
\end{proof}

\noindent Note that this proof only
uses the strong factorizing cactus
property and the concentration of the
traffic distribution, which explains
why~\Cref{thm:orthogonal-amp} also holds
for non-orthogonally invariant matrix
models such as Wigner matrices (\Cref{sec:traffic-wigner}).

\subsubsection{Punctured random and deterministic matrices}
\label{sec:dynamics-reg}

The punctured matrices studied in~\Cref{sec:universality} do not satisfy the strong cactus property, so we cannot directly apply~\Cref{thm:full-onsager} to derive an AMP iteration for them. However, a reduction allows us to derive the state evolution of punctured orthogonally invariant random matrices from that of their unpunctured counterparts. These matrices are central because, by~\Cref{thm:universality-new}, they provide an intermediate step in deriving the state evolution of sequences of punctured {\em deterministic} matrices satisfying~\Cref{assumption:punctured}.

Note that a GFOM run on a punctured matrix must be initialized with a random vector $\bmx_0 \sim \calN(\bm{0}, \bI)$, rather than $\bm{x}_0 = \bm{1}$, to avoid triviality.

\begin{theorem}[State evolution for punctured matrices]\label{prop:hadamard-explicit-state}
    Let $\bH = \bH^{(n)} \in \R_\sym^{n \times n}$ be a sequence of orthogonally invariant random matrices satisfying~\cref{eq:tightnessNorm} and converging in tracial moments in $L^2$ to a probability measure with free cumulants $(\kappa_q)_{q\ge 1}$. 
    Let $\bA$ denote the puncturing of $\bH$ (\Cref{def:puncturing}).
    Let $f_t : \R \to \R$ be polynomial functions with $f_0(x)=x$, and consider the pGFOM:
    \begin{alignat*}{2}
        \bm x_0 &\sim \calN(\bm 0, \bm I)\,,\quad &\bm x_t &= \bm A { f_{t-1}(\bm x_{t-1})} - \sum_{s=0}^{t-1} \kappa_{t-s} \left(\prod_{r=s+1}^{t-1} \langle f_r'(\bm x_r)\rangle\right) (f_s(\bm x_s) - \langle f_s(\bm x_s)\rangle \bm{1})\quad \forall t\ge 1\,.\label{eq:ampPunctured}
    \end{alignat*}
    Then for any $t\ge 1$ and any polynomial $\varphi:\R^t \to \R$, we have
    \[
        \lim_{n\to\infty}  \E_{\bm H,\bm x_0} \langle \varphi(\bm x_1, \ldots, \bm x_t)\rangle = \E \varphi(X_1, \ldots, X_t)\,,
    \]
    where $(X_t)_{t\ge 1}$ is a centered Gaussian process with covariance given by
    \begin{align*}
        \E\left[X_s X_t\right] &= \sum_{s' = 0}^{s-1}\sum_{t'=0}^{t-1} \kappa_{s-s'+t-t'} \left(\prod_{r= s'+1}^{s-1} \E f'_r(X_r) \right) \left(\prod_{r=t'+1}^{t-1} \E f'_r(X_r)\right) \E \left[\overline{F_{s'}} \, \overline{F_{t'}}\right]\quad \forall s,t\ge 1\,,\\
        \overline{F_0} &\defeq 1\,,\quad \overline{F_t} \defeq f_t(X_t) - \E f_t(X_t)\quad \forall t\ge 1\,.
    \end{align*}
\end{theorem}

\noindent By~\Cref{thm:universality-orthogonal}, the conclusion of~\Cref{prop:hadamard-explicit-state} also holds for any sequence of deterministic matrices satisfying the delocalization assumption~\Cref{assumption:punctured} and having a limiting diagonal distribution that factorizes over cycles (that is, matches the diagonal distribution of some orthogonally invariant random matrix ensemble).
In particular, the conclusion holds for the Walsh-Hadamard matrices and the Discrete Cosine and Sine Transform matrices, for which the $\kappa_q$ are the free cumulants of the \rom  (\cref{eq:freeCumulantROM}).

The proof of~\Cref{prop:hadamard-explicit-state} proceeds by reducing to the following iteration on the original, non-punctured matrix, initialized at the all-ones vector:
\begin{alignat}{2}
    \bm u_0 &= \bm 1\,,\qquad
    \bm u_t
    = \bm H \overline{\bm f_{t-1}}
    - \sum_{s=0}^{t-1} \bm b_{s,t} \cdot \overline{\bm f_s} \quad \forall t\ge 1\,,\nonumber\\
    \text{where } {\bm b}_{s,t}[i]
    &\defeq \sum_{\substack{i_s,\ldots,i_{t-1}\in[n] \textnormal{ distinct}\\i_s=i}}
    \left(\prod_{r=s+1}^{t-1} \bm H[i_{r-1}, i_{r}] \bm f'_{r}[i_{r}] \right)
    \bH[i_{t-1},i_s]\quad \forall t> s\ge 0\,,\label{eq:backtrackPunctured}\\
    \overline{\bm f_0} &\defeq \bm u_0\,,\quad \overline{\bm f_t} \defeq \bm \Pi f_t(\bm u_t)\quad \forall t\ge 1\,.\nonumber
\end{alignat}

\begin{lemma}\label{lem:switchInit}
    For any $t\ge 1$ and any polynomial $\varphi:\R^t \to \R$,
    \begin{align*}
        \lim_{n\to\infty} \E_{\bm H,\bm x_0} \langle \varphi(\bm x_1, \ldots, \bm x_t)\rangle &= \lim_{n\to\infty} \E_{\bm H} \langle \varphi(\bm u_1, \ldots, \bm u_t)\rangle\,.
    \end{align*}
\end{lemma}

\noindent The proof of~\Cref{lem:switchInit}
is deferred to~\Cref{appendix:initialization}.

\begin{proof}[Proof of~\Cref{prop:hadamard-explicit-state}]
    We apply~\Cref{thm:orthogonal-amp} to $\bm u_t$ after replacing iteratively each occurrence of $\bm \Pi f_t(\bm u_t)$ by $f_t(\bm u_t) - \E\left[f_t(U_t)\right] \cdot \bm 1$ (where $U_t$ is the asymptotic state of $\bm u_t$ as predicted by~\cref{thm:orthogonal-amp}). By~\Cref{lem:asymptotic-state-empirical-average}, this transformation does not change the asymptotic state of $\bm u_t$.
    The state evolution formula for polynomial test functions then transfers to $\bm x_t$ by~\Cref{lem:switchInit}.
\end{proof}

\subsubsection{Block-structured random matrices}

\newcommand{\Community}{\textnormal{Community}}

Our final example is the class of block-structured matrices whose blocks satisfy the factorizing strong cactus property, which we introduced in~\Cref{sec:wigner-block}.
As anticipated in~\Cref{ex:block-structured}, these matrices do not
themselves
satisfy the factorizing strong
cactus property.
Therefore, we start by describing the random
limit $Z^\infty_{\calC_1}$.


\begin{lemma}\label{lem:block-matrix-cactus-limits}
    Let $q \in \N$. For $r,c\in [q]$, let $\bA_{r,c} = \bA_{r,c}^{(n)} \in \R^{\frac nq \times \frac nq}_\sym$ be a sequence of symmetric random matrices such that $\bm A_{r,c} = \bm A_{c,r}$.
    Let $\bA \in \R^{n \times n}_\sym$ be the block matrix with blocks $(\bA_{r,c})_{r,c \in [q]}$.
    Assume that each $\bm A_{r,c}$ satisfies the factorizing strong cactus property,
    and $(\bm A_{r,c})_{1 \leq r\leq c \leq q}$
    are asymptotically traffic independent. Let $(\kappa^{\{r,c\}}_\ell)_{\ell\ge 1}$ be the limiting
    free cumulants of $\bm A_{r,c}$.
    Then,
    \[
        (\block(i), \bmz_{\cC_1}(\bm A)[i]) \tod  (R, Z_{\cC_1}^\infty(R))\,, \qquad i \sim \Unif([n])\,, \quad R \sim \Unif([q])\,,
    \]
    where  the (deterministic) sequence $Z_{\cC_1}^\infty(r)$ for $r\in [q]$ is defined recursively by:
    \begin{enumerate}[(i)]
        \item For the singleton cactus, $Z^\infty_{\mathrm{singleton}}(r) \defeq 1$.
        
        \item Suppose $\sigma \in \mathcal C_1$ is rooted at a vertex $u_1$ of degree $2$. Let $(u_1, \ldots, u_{\ell})$ be
        the cycle incident to the root.
        Let
        $\sigma_2, \dots, \sigma_{\ell} \in \mathcal C_1$ be the rooted cactuses
        attached to the vertices of the cycle. Then
        \[
            Z^\infty_\sigma(r) \defeq
            \begin{cases}
                \displaystyle
                \sum_{c\in [q]} \left[\kappa_\ell^{\{r,c\}}
                \prod_{\substack{k=2\\k\textnormal{ odd}}}^{\ell} Z^\infty_{\sigma_k}(r) \prod_{\substack{k=2\\k\textnormal{ even}}}^{\ell} Z^\infty_{\sigma_k}(c)\right]
                & \text{if $\ell$ is even} \\
                \displaystyle
                \kappa_\ell^{\{r,r\}}
                \prod_{k=2}^{\ell} Z^\infty_{\sigma_k}(r)
                & \text{if $\ell$ is odd}
            \end{cases}
        \]
    
        \item If $\sigma \in \mathcal C_1$ decomposes as
        $\sigma = \bigoplus_{k=1}^\ell \sigma_k$, then $Z^\infty_\sigma(r) \defeq \prod_{k=1}^\ell Z^\infty_{\sigma_k}(r)$.
    \end{enumerate}
    In particular, the law of the limit $Z_{\calC_1}^{\infty}$ is $\Unif(\{Z_{\calC_1}^{\infty}(r): r \in [q]\})$.
\end{lemma}

\noindent The proof is deferred to~\Cref{app:block-matrix}. By unicity of the limit in distribution,~\Cref{lem:block-matrix-cactus-limits}, together with~\Cref{thm:convergence-distribution}, determines
the law of $Z^\infty_{\calA_1}$.
The joint convergence with $\block(i)$ clarifies the source of the randomness
of $Z^\infty_{\calC_1}$: it arises because of the random choice of the block an entry belongs to in the $\samp(\cdot)$ operation.

Using~\Cref{lem:block-matrix-cactus-limits},
we can specialize the treelike
AMP iteration and its state evolution
to concrete block-structured models.
We start with block GOE matrices (\cref{def:blockgoe}).
A family of AMP iterations for such matrices was derived in~\cite{rangan2011generalized, javanmard2013state}. As we will discuss below, these iterations
have the same asymptotic state as
treelike AMP.

\begin{theorem}[State evolution for the block GOE model]\label{thm:block-goe-state}
    Let $\bA \sim \bGOE(n, \bm\Sigma)$, where  $\bm \Sig \in \R_{\ge 0}^{q \times q}$ is a symmetric matrix.
    Given polynomial functions $f_t : \R \to \R$, let
    \begin{alignat}{2}
        \bmx_0 &= \bm{1}\,,\quad \bm x_1 = \bm A f_0(\bm x_0)\,, \quad &\bmx_t &= \bA f_{t-1}(\bm x_{t-1}) - (\bm A^{\odot 2}  f_{t-1}'(\bm x_{t-1})) \cdot f_{t-2}(\bm x_{t-2})\quad \forall t\ge 2\,.\label{eq:amp-block-goe}
    \end{alignat}
    Then the asymptotic state of $(\bm x_t)_{t\ge 1}$
    is a mixture $\frac 1q\sum_{r \in [q]} \mu_r$ where $\mu_r$ denotes the law of a centered Gaussian process $(X_t)_{t \geq 1}$ with covariance kernel $\bm \Gamma_r$ defined recursively by
    \[
        \bm \Gamma_{r}[s,t] = \sum_{c\in  [q]}\bigg[\bSig[r,c] \E_{(X_{T})_{T \ge 1} \sim \mu_c} \left[f_{s-1}(X_{s-1})f_{t-1}(X_{t-1})\right]\bigg]\quad \forall r \in [q]\,,\quad  \forall s,t\ge 1\,,
    \]
    with $X_0\defeq 1$.
\end{theorem}

\begin{proof}
    By the discussion after~\Cref{thm:male-traffic-independence}, $\bA$ has a traffic distribution and satisfies the strong cactus
    property\footnote{One can also verify that it satisfies~\cref{eq:tightnessNorm}. But note that this was only needed in the proof of~\Cref{thm:convergence-distribution} to ensure the existence of the limit $Z^\infty_{\calC_1}$, which we established directly in~\Cref{lem:block-matrix-cactus-limits}.}, so it satisfies
    the assumption of~\Cref{thm:full-onsager}.
    We consider the treelike AMP iteration $\bm x_t$
    in~\cref{eq:tree-amp} applied to
    $\bm A$.
    We show that this iteration
    has the same asymptotic state
    as~\cref{eq:amp-block-goe} by
    simplifying the Onsager correction term.
    
    The free cumulants of the GOE are $0$ except for $\kappa_2$, so by~\cref{lem:block-matrix-cactus-limits}, the only asymptotically non-negligible cactuses are those such that every cycle is a 2-cycle.
    For any
    $s < t-2$, $\bm b_{s,t}$ contains an injective cycle
    of length larger than $2$ that cannot be
    destroyed by later operations. For $s=t-2$,
    we have:
    \begin{align*}
        \bmb_{t-2,t}[i] &= \sum_{\substack{j = 1\\j \neq i}}^n \bA[i,j]^2 \bmf'_{t-1}[j] = (\bm A^{\odot 2} \bm f_{t-1}')[i] - \bm A[i,i]^2 \bm f_{t-1}'[i]\,.
    \end{align*}
    Both $\bm A[i,i]^2 \bm f_{t-1}'[i]$ and
    $\bm b_{t-1,t}$ contain a self-loop that also
    cannot be destroyed by later operations.
    In conclusion, the treelike AMP algorithm from~\cref{eq:tree-amp} and the iteration in~\cref{eq:amp-block-goe} are equal up to negligible diagrams.
    By~\Cref{thm:full-onsager},
    the asymptotic state $(X_t)_{t\ge 1}$ of $(\bmx_t)_{t\ge 1}$ in~\cref{eq:amp-block-goe} exists and
    is Gaussian conditionally on
    $Z^\infty_{\calC_1}$, and so, in the construction from~\Cref{lem:block-matrix-cactus-limits}, it is Gaussian conditionally on the random variable $R$. 

    Next, we specialize the covariance formula given by~\Cref{lem:cov-cactus}. Since
    only cactuses of 2-cycles are nonzero in the traffic distribution of $\bm A$ (this may be induced from~\Cref{lem:block-matrix-cactus-limits}), only the term
    for $s'=s-1$ and $t'=t-1$ is non-negligible
    in the expansion of $\frac 1 n \E \bm x_s \cdot \bm x_t$ given by~\Cref{lem:cov-cactus}. The expansion into cactuses of that
    term is obtained by grafting together a 2-cycle at the root, and cactuses of 2-cycles
    from $\bm f_{s-1}$ and $\bm f_{t-1}$ at the
    child of the root. 
    Applying the recursive formula for $Z^\infty_{\calC_1}(r)$ in~\Cref{lem:block-matrix-cactus-limits}, we obtain:
    \begin{align*}
        \E \left[X_s X_t\mid R=r\right]
        &=  \sum_{c\in [q]} \bm \Sigma[r,c] \E\left[f_{s-1}(X_{s-1}) f_{t-1}(X_{t-1})\mid R=c\right]\qquad \forall r\in [q]\,.
    \end{align*}
    Thus, we have shown that, conditionally on $R\sim \Unif([q])$, $(X_t)_{t\ge 1}$
    is a Gaussian process with the required covariance. The result follows by taking $\mu_r$ to be the law of $(X_t)_{t\ge 1}$ conditionally
    on $R=r$.
\end{proof}

To illustrate the modularity of
our approach, we also study a different
block-structured matrix
model whose blocks
are not all GOE.

\begin{theorem}[State evolution for the community model]\label{thm:communitySimple}
    Let $\bm M\in \R_{\sym}^{\frac n q\times \frac n q}$ be
    an orthogonally invariant
    random matrix converging in
    tracial moments to a probability measure with free cumulants $(\kappa_q)_{q\ge 1}$
    such that $\kappa_2 = \frac 1 q$.
    Let $\bm A$ be the random symmetric $n\times n$ matrix with blocks $(\bA_{r,c})_{r,c \in [q]}$ given by $\bA_{1,1} = \bM$ and 
    for all $1\le r\le c\le q$ with $(r,c)\neq (1,1)$, 
    $\bA_{r,c}$ are i.i.d. $\frac n q\times \frac n q$ GOE matrices with entries of variance $\frac 1n$ (and we set $\bm A_{r,c} =\bm A_{c,r}$).
    
    Let $\bm x_t$ be the treelike AMP iteration~\cref{eq:tree-amp} run
    on $\bm A$ with arbitrary polynomial
    nonlinearities.
    Then the asymptotic state $(X_t)_{t\ge 1}$ of $(\bm x_t)_{t\ge 1}$ is the mixture $(1-\frac 1q)\mu_0 + \frac 1 q \mu_1$, where $\mu_i$ is the law
    of a centered Gaussian process
    $(X_t)_{t\ge 1}$ with covariance kernel $\bm \Gamma_i$
    defined recursively by, for
    all $s,t\ge 1$:
    \begin{align*}
        & \bm \Gamma_0[s,t] = \E\left[F_{s-1}F_{t-1}\right]\,\\
        & \bm \Gam_1[s,t] = \E\left[F_{s-1}F_{t-1}\right]  + \underset{(s',t') \neq (s-1,t-1)}{\sum_{s'=0}^{s-1} \sum_{\substack{t'=0}}^{t-1}} \kappa_{s-s'+t-t'} 
        \left(\prod_{r=s'+1}^{s-1} \E_{\mu_1}F'_r\right) \left(\prod_{r=t'+1}^{t-1}\E_{\mu_1} F'_r \right) \E_{\mu_1}\left[F_{s'}F_{t'}\right]\,,\\
        & F_t := f_t(X_t), \qquad F'_t := f'_t(X_t), \qquad X_0 = 1\,,
    \end{align*}
    where $\E_{\mu_1}$ denotes
    expectation with respect to $(X_t)_{t\ge 1}\sim \mu_1$.
\end{theorem}

\begin{proof}
    The assumptions of~\Cref{lem:block-matrix-cactus-limits,thm:communitySimple} are satisfied. All
    blocks except the one in position $(1,1)$ have the same
    free cumulants (the GOE free cumulants, normalized so that $\kappa_2 = \frac 1 q$). Therefore, in the construction of~\Cref{lem:block-matrix-cactus-limits}, we have $Z^\infty_{\calC_1}(r) = Z^\infty_{\calC_1}(s)$ for all $r,s>1$. Let $\mu_0$ (resp. $\mu_1$) be
    the law of the asymptotic state
    $(X_t)_{t\ge 1}$ of the treelike AMP iteration $(\bm x_t)_{t\ge 1}$ conditioned on $R>1$ (resp. $R=1$). By~\Cref{thm:full-onsager}, both $\mu_0$ and $\mu_1$ are the laws of centered Gaussian processes.
    It remains to specialize the formula of \Cref{lem:cov-cactus} for their covariance 
    to the present setting. 
    
    Conditionally
    on $R>1$ (that is, outside the community), only 2-cycles
    at the root contribute to $Z^\infty_{\calC_1}$. Thus, by
    combining the strong cactus property,~\Cref{lem:cov-cactus},
    and~\Cref{lem:block-matrix-cactus-limits}, we obtain
    \begin{align*}
        \E \left[X_s X_t\mid R >1\right] &= \frac 1q \left(\E \left[f_{s-1}(X_s) f_{t-1}(X_t) \mid R = 1\right] + (q-1) \E \left[f_{s-1}(X_s) f_{t-1}(X_t) \mid R > 1\right]\right)\\
        &= \E \left[f_{s-1}(X_s) f_{t-1}(X_t)\right]\,.
    \end{align*}
    
    Conditionally on $R=1$ (that is, inside
    the community), we also obtain 
    a contribution of $\E \left[f_{s-1}(X_s) f_{t-1}(X_t)\right]$ from the term
    $s'=s-1$ and $t'=t-1$ in~\Cref{lem:cov-cactus} (again using the normalization $\kappa_2 = \frac 1q$ inside the community).
    For all of the remaining terms $s',t'$,
    when $s-s' + t-t'$ is an even
    integer larger than $2$, we obtain
    a contribution only from $c=1$
    in~\Cref{lem:block-matrix-cactus-limits}, namely
    \[
        \kappa_{s-s'+t-t'} \E\left[ F_{s'} F_{t'}\mid R=1\right]\prod_{r=s'+1}^{s-1} \E\left[ F'_{r}\mid R=1\right]\prod_{r=t'+1}^{t-1} \E\left[ F'_{r}\mid R=1\right]\,.
    \]
    When $s-s'+t-t'$ is odd,~\Cref{lem:block-matrix-cactus-limits} yields exactly the same expression as the even case.
    Altogether, we obtain the recursion
    \begin{align*}
        &\qquad\E\left[X_s X_t\mid R=1\right] = \E \left[F_{s-1} F_{t-1}\right] + \\
        &\underset{(s',t') \neq (s-1,t-1)}{\sum_{s'=0}^{s-1} \sum_{\substack{t'=0}}^{t-1}} \kappa_{s-s'+t-t'} \E\left[ F_{s'} F_{t'}\mid R=1\right] \prod_{r=s'+1}^{s-1} \E\left[ F'_{r}\mid R=1\right]\prod_{r=t'+1}^{t-1} \E\left[ F'_{r}\mid R=1\right] \,.
    \end{align*}
    These are the desired covariance
    formulas for $\mu_0$ and $\mu_1$, and the mixing weights
    of the events $(R=1)$ and $(R>1)$ are indeed $\frac 1q$ and $1-\frac 1 q$, respectively.
\end{proof}

\subsubsection{Further extensions}

There are several possible technical extensions of the methods we have developed here,
whose full development is left for future work.

First,~\Cref{lem:block-matrix-cactus-limits} applies to general orthogonally invariant distributions within the blocks, not just the GOE. In principle, one can then derive a corresponding state evolution formula mechanically for non-identically distributed orthogonally invariant blocks with arbitrary free cumulants, although the resulting expression is quite complicated.

Second, for technical reasons, we assumed that the blocks are square and symmetric, so that we could work with undirected graphs. The results of~\cite{male2020traffic,cebron2024traffic} extend to general matrices, and our techniques should also extend to the setting of varying block sizes and asymmetric matrices, leading to non-uniform mixtures in the recursion for the covariance kernel.

One caveat of the treelike AMP algorithm
is that the Onsager correction term
in~\cref{eq:tree-amp} is not obviously efficient
to compute in practice.\footnote{The $\bm b_{s,t}$ can be approximated with high probability to negligible error for all
$0\le s<t\le T$ in time $2^{O(T)}\poly(n)$ using the color coding
technique~\cite{colorCoding, heavyTailedWigner}, but the exponential
dependence on $T$ makes this
algorithm impractical to implement
for large $T$.}
On the other hand,
the vectors $\bb_{s,t}$ have asymptotically constant entries in many settings,
so that the Onsager correction can be replaced by a simpler asymptotically equivalent term, like in \Cref{thm:orthogonal-amp,prop:hadamard-explicit-state}.
This should also hold for block-structured models, as in the generalized AMP algorithm of Javanmard and Montanari~\cite{javanmard2013state}.
For example, \cref{eq:amp-block-goe} is expected to be asymptotically equivalent to:
\begin{alignat*}{2}
    \bmx_0 &= \bm{1}\,, \qquad &\bmx_t &= \bA f_{t-1}(\bm x_{t-1}) - \bmb_{t-2,t} \cdot f_{t-2}(\bm x_{t-2})\,,\\
     && \bmb_{t-2,t}[i] &= \sum_{c = 1}^q \matSig[\block(i), c] \langle f'(x_{t-1}) \cdot \bm{1}_{\block = c} \rangle\,,
\end{alignat*}
where $\bm{1}_{\block = c} \in \{0,1\}^{n}$ indicates the entries in block $c \in [q]$.
The treelike AMP algorithm for~\cref{thm:communitySimple} is expected to be asymptotically equivalent to:
\begin{alignat*}{2}
    \bmx_0 &= \bm{1}\,, \qquad &\bmx_t &= \bA \bmf_{t-1} - \langle \bmf'_{t-1} \rangle \bmf_{t-2}
    - \sum_{\substack{s = 0\\ s \neq t-2}}^{t-1} \kappa_{t-s} \left(\prod_{r=s+1}^{t-1} \langle \bmf_r' \cdot \bm{1}_{\block = 1} \rangle \right) \bmf_s \cdot \bm{1}_{\block = 1}\,, \\
    \bmf_{t} &\defeq f_t(\bmx_t)\,, \qquad &\bmf'_t &\defeq f'_t(\bx_t)\,,
\end{alignat*}
where $\bm{1}_{\block = 1} \in \{0,1\}^n$ indicates the entries in the first block.
Because these expressions involve the blockwise indicators
$\bm{1}_{\block = c}$, they could be represented and analyzed using an extended
diagram basis in which certain indices are constrained to lie in a prescribed
block. We leave the full development of this extension to future work.

A final open question is to characterize traffic distributions satisfying the
(not necessarily factorizing) strong cactus property. Sequences of block matrices with orthogonally invariant
blocks provide one general construction of matrices with the strong cactus
property. If a sequence of matrices has the strong cactus property, must its
traffic distribution arise as the limit (in an appropriate sense) of traffic
distributions of block matrices with orthogonally invariant blocks (allowing the number of blocks to tend to infinity)?